Cancer is a mixed group of diseases that can be hematologically (blood) related or 
tumour related \cite{NCI}. 
They are all characterised by their uncontrolled cellular growth that is triggered by 
genetic mutations and epigenetic changes (environmental factors) alterations 
\cite{NCI}. Traditionally, treatment of cancer relied on 
generalised treatment protocols that were based on the anatomical site of the cancer. 
Modern protocols have shifted to what can be termed as precision medicine, which is 
aimed at providing customised individual patient care. Precision pathology in particular, 
aims to tailor treatment to patients based on the individual molecular profile of the 
patient’s cancer. \cite{Klauschen2024}. Achieving this requires the integration of 
high-dimensional genomic data, such as gene expression profiles (transcriptomics), 
to predict how a specific cancer will respond to targeted therapies \cite{Klauschen2024}.
\\\\
Precision pathology has evolved from manual microscopic inspection methods to AI-driven 
diagnostics. Machine learning models are now capable of analysing complex biological 
datasets to predict patient survival outcomes and drug sensitivities \cite{Kovalev2021}. 
A particular case is that of machine learning survival models. Machine learning survival 
models are predictive models designed to estimate the time-to-event outcomes, for example 
how long a patient is likely to live after diagnosis and or treatment \cite{Kovalev2021}. 
This is useful in determining the potency of a treatment plan based on the patient's characteristics 
\cite{Kovalev2021}. 

% [Paste Section 3.2.1 Genomics and the Interpretability Gap here]
\subsubsection{Genomics and the Interpretability Gap}
The precision of machine learning models in precision medicine is driven by massive datasets 
such as the Genomics of Drug Sensitivity in Cancer (GDSC). In genomic-based cancer research, 
datasets like (GDSC) provide thousands of gene expression variables \cite{Baptista2021,Yang2013}. 
The GDSC dataset is a large publicly available database that is designed to help researchers 
discover therapeutic biomarkers (e.g. the expression levels of the EGFR gene) and predict 
how specific tumours will respond to various cancer treatments \textcolor{red}{(revisit intro of paper)}
\cite{Yang2013}. It is practically, as well as financially, infeasible to test every 
candidate drug directly on human patients. Therefore, cancer cell lines are used as a proxy 
for real tumours \cite{Yang2013}. Cancer cell lines are cells that have been extracted 
and grown in a laboratory and still retain the genetic and molecular characteristics of the 
cancers they originate from \cite{Yang2013}. Drug response is measured in the form of IC50 
values \cite{Yang2013}.
\\\\
While DNNs excel at finding patterns in high-dimensional datasets such as GDSC, they largely 
operate as black boxes \cite{Partin2023,Klauschen2024}. To elaborate, current research 
has successfully used AI to identify biomarkers that correlate with treatment resistance or treatment 
efficacy \cite{Partin2023,Klauschen2024,Baptista2021}. Despite this success in 
predictive accuracy, a significant trust gap exists, which is that clinicians are not able to easily 
verify whether a model’s prediction is based on genuine oncogenic pathways or statistical noise 
\cite{Klauschen2024}. \cite{Klauschen2024} argue that for AI to be integrated into clinical 
workflows, it must provide "explainable" outputs that align with established medical knowledge.
\\\\
In XAI, unreliability of explanations is related to data corruption or noise \cite{Klauschen2024}. 
In cancer genomics, small variations in the sample preparation can lead to batch effects which AI models 
learn as shortcuts to high accuracy or Clever Hans \cite{Partin2023,Klauschen2024}. Batch 
effects are non-biological, technical variations that are introduced into the data \cite{Yu2024, Goh2022}. 
Batch effects are not related to the variables of interest in the study and come about due to data being 
collected from different labs or machines, different assessment pipelines or differing experimental 
conditions over time \cite{Yu2024}. 
\\\\
Survival analysis research machine learning models shows that black box models can produce highly 
accurate predictions scores while hiding the underlying counterfactuals - specific changes in a patient's 
genetic expression can alter the clinical outcome or decision \cite{Kovalev2021}. Furthermore, if an 
XAI method identifies a "rule" for a cancer diagnosis that is mathematically sound but biologically 
incoherent (e.g., incorrectly linking a housekeeping gene to tumour suppression), it can lead to dangerous 
medical misinterpretations \cite{Partin2023,Klauschen2024}. This underscores the need for the 
rigorous "Audit and Verify" proposed in this research project to ensure that interpretability serves as a 
safeguard for clinical truth rather than just a justification for model output. 
